\begin{enumerate}[label=\textbf{\arabic*.}]

    \item \underline{\textbf{Spolupráca s vedúcim}} \hfill \underline{\textbf{max. 20}}
    
    \begin{enumerate}
        \item \textit{Efektivita konzultácií} \hfill \textit{max. 15}
        \begin{itemize}
            \item \textbf{Do 15b} – Študent priebežne konzultoval postup a výsledky práce s vedúcim, riadil sa pokynmi vedúceho a vykonal korekcie ku všetkým pripomienkam.
            \item \textbf{Do 11b} – Študent priebežne konzultoval postup a výsledky práce s vedúcim, ale viaceré pokyny alebo pripomienky nerešpektoval.
            \item \textbf{Do 7b} – Študent priebežne konzultoval postup a výsledky, ale neriadil sa pokynmi vedúceho a výsledky prevažne nezodpovedajú predstave vedúceho.
            \item \textbf{0b} – Študent nekonzultoval postup ani výsledky a tie nezodpovedajú predstave vedúceho.
        \end{itemize}

        \item \textit{Iniciatíva pri získavaní podkladov} \hfill \textit{max. 5}
        \begin{itemize}
            \item \textbf{Do 5b} – Študent využil okrem odporúčanej literatúry aj iné literárne zdroje a publikácie (knihy, odborné časopisy, vedecké články).
            \item \textbf{Do 3b} – Študent využil odporúčanú literatúru a elektronické zdroje.
            \item \textbf{Do 1b} – Študent nedostatočne využil odporúčanú literatúru a elektronické zdroje.
        \end{itemize}
    \end{enumerate}

    \item \underline{\textbf{Splnenie cieľov zadania}} \hfill \underline{\textbf{max. 20}}
    
    \begin{enumerate}
        \item \textit{Kvantitatívne plnenie úloh} \hfill \textit{max. 10}
        \begin{itemize}
            \item \textbf{Do 10b} – Všetky úlohy zadania boli riešené v plnom rozsahu.
            \item \textbf{Do 5b} – Niektorá z úloh si vyžadovala viac priestoru než jej študent venoval.
        \end{itemize}

        \item \textit{Kvalitatívne plnenie úloh} \hfill \textit{max. 10}
        \begin{itemize}
            \item \textbf{Do 10b} – Výsledky realizovaných úloh úplne spĺňajú očakávania vedúceho.
            \item \textbf{Do 7b} – Niektorá z úloh, bez ohľadu na rozsah vykonanej práce, nemá výstupy v požadovanej kvalite.
            \item \textbf{0b} – Prevažná časť výstupov realizovaných úloh nespĺňa očakávania vedúceho.
        \end{itemize}
    \end{enumerate}

    \item \underline{\textbf{Kvalita priebežnej práce}} \hfill \underline{\textbf{max. 30}}
    
    \begin{enumerate}
        \item \textit{Organizácia práce} \hfill \textit{max. 10}
        \begin{itemize}
            \item \textbf{Do 10b} – Študent na úlohách pracoval pravidelne, rovnomerne si rozložil objem vykonávanej práce. Počas celej doby riešenia práce nikdy nestagnoval a vždy bolo možné pozorovať progres.
            \item \textbf{Do 7b} – Študent na úlohách pracoval nepravidelne alebo nárazovo a občas stagnoval. Nemalo to však výrazný dopad na kvalitu vykonávanej práce.
            \item \textbf{Do 4b} – Študent na úlohách pracoval nepravidelne alebo nárazovo a výrazne stagnoval, čo sa prejavilo na kvalite vykonávanej práce.
        \end{itemize}

        \item \textit{Efektivita práce} \hfill \textit{max. 20}
        \begin{itemize}
            \item \textbf{Do 20b} – Študent efektívne využíval a vhodne prepájal potrebné metódy a nástroje, a teda maximum svojho úsilia mohol vynaložiť na kvalitatívne a kvantitatívne plnenie úloh.
            \item \textbf{Do 15b} – Študent venoval väčšinu svojho úsilia na zvládanie potrebných podporných prostriedkov alebo na vykonávanie okrajových činností.
            \item \textbf{Do 10b} – Študent venoval väčšinu svojho úsilia na zvládanie okrajových činností alebo na tvorbu podporných systémov. Plneniu podstaty úloh zadania venoval len minimum priestoru.
        \end{itemize}
    \end{enumerate}

    \item \underline{\textbf{Dodržanie časového harmonogramu}} \hfill \underline{\textbf{max. 10}}
    
    \begin{itemize}
        \item \textbf{Do 10b} – Študent dodržoval dohodnuté termíny, čiastkové výstupy prezentoval vedúcemu podľa dohody a prácu odovzdal načas.
        \item \textbf{Do 6b} – Študent prevažne dodržoval dohodnuté termíny. Pracoval s časovým sklzom, no konečnú prácu odovzdal načas.
        \item \textbf{Do 3b} – Študent pracoval s výrazným časovým sklzom a prácu odovzdal až počas predĺženého termínu.
    \end{itemize}

    \item \underline{\textbf{Samostatný prístup k práci}} \hfill \underline{\textbf{max. 20}}
    
    \begin{itemize}
        \item \textbf{Do 20b} – Študent bol iniciatívny, sám navrhoval postupy práce a riešenia, ktoré zodpovedali predstave vedúceho. Všetky potrebné poznatky si samostatne doštudoval.
        \item \textbf{Do 12b} – Študent iniciatívne pristupoval k riešeniu práce, no nenavrhoval vhodné postupy a riešenia. Spoliehal sa prevažne na pokyny vedúceho.
        \item \textbf{Do 8b} – Študent sa pri riešení práce iniciatívne neprejavoval. Pracoval len podľa pokynov vedúceho.
    \end{itemize}

\end{enumerate}